\begin{textsample}{POS Dim 3 – human – Score 7.00 – mg/mg\_ef\_linguagens\_lp\_1\_p1.txt}  \label{ex:f3_pos_007}
O componente \textbf{curricular} Língua Portuguesa fundamenta -se nos Parâmetros \textbf{Curriculares} Nacionais ( PCNs ), nas \textbf{Diretrizes} \textbf{Curriculares} Nacionais ( DCN ), na Base \textbf{Nacional} Comum \textbf{Curricular} ( BNCC ).
Todo o trabalho apresentado para este componente \textbf{Curricular}, no Currículo Referência de Minas Gerais, está de \textbf{acordo} com os pressupostos teóricos-metodológicos apresentados pela BNCC, não havendo exclusões devido ao seu caráter normativo.
O \textbf{documento} foi complementado com as orientações produzidas pelo \textbf{Estado} de Minas Gerais, acrescidas de atualizações e \textbf{especificidades} que agora se fazem necessárias à criação do novo currículo, de forma que atenda às demandas surgidas pelo regime de colaboração.

% matched lemmas: acordo, curricular, diretriz, documento, especificidade, estado, nacional
\end{textsample}
